\documentclass[12pt,a4paper]{article}
\usepackage[utf8]{inputenc}
\usepackage[margin=1in]{geometry}
\usepackage{amsmath,amssymb}
\usepackage{enumitem}
\usepackage{xcolor}
\usepackage{hyperref}
\usepackage{fancyvrb}
\usepackage{fancyhdr}
\hypersetup{colorlinks=true,linkcolor=black,urlcolor=blue}

\setlist{nosep,topsep=2pt,parsep=2pt,itemsep=2pt}

\newcommand{\reviewercomment}[1]{\par\smallskip\noindent\textit{Reviewer:}\par\smallskip\begingroup\itshape #1\par\endgroup\smallskip}
\newcommand{\ourresponse}[1]{\par\noindent\textbf{Response:} #1\par\medskip}

% Footer: contact email on every page
\pagestyle{fancy}
\fancyhf{}
\renewcommand{\headrulewidth}{0pt}
\fancyfoot[L]{\small Correspondence: \href{mailto:francesca.castaldo@bcom.one}{francesca.castaldo@bcom.one}; \href{mailto:francesca.castaldo@neuroelectrics.com}{francesca.castaldo@neuroelectrics.com}}
\fancyfoot[R]{\small \thepage}
\fancypagestyle{plain}{%
  \fancyhf{}%
  \renewcommand{\headrulewidth}{0pt}%
  \fancyfoot[L]{\small Correspondence: \href{mailto:francesca.castaldo@bcom.one}{francesca.castaldo@bcom.one}; \href{mailto:francesca.castaldo@neuroelectrics.com}{francesca.castaldo@neuroelectrics.com}}%
  \fancyfoot[R]{\small \thepage}%
}

\title{Response to Reviewer\\[0.4em]
\large \emph{The Rosetta Stone of Neural Mass Models}\\[0.2em]
\normalsize Francesca Castaldo, Raul de Palma Aristides, Pau Clusella,\\
Jordi Garcia-Ojalvo, Giulio Ruffini\\[0.3em]
\normalsize Submitted to \textit{Physics Reports}}
\author{}
\date{1 May 2026}

\begin{document}
\maketitle

We are grateful for the reviewer's comments, which have improved the framing, lineage, and technical accuracy of the manuscript. We respond point by point below; the reviewer's comments are reproduced verbatim, and revised text in the manuscript is wrapped in \Verb|\rev{...}| (rendered in blue) to make the changes locatable.

\section*{Overview}

\reviewercomment{Neural mass models are now a mainstay of large-scale brain modelling, and have been a topic of theoretical development with applications to brain science, including cognition and neuroimaging, for some years. A review that synthesises knowledge and advances in this space for a physics audience, allowing new researchers to survey the field and existing ones to keep up to date, is a good topic for Physics Reports. As it currently stands I feel that this review could be strengthened in a number of ways. Currently it is not sufficiently informative or slightly misleading / misrepresentative on a number of counts. I appreciate the challenge of writing such a review (to be intelligible to non-specialists within physics, so readers outside a narrow subfield can still grasp the main ideas). At the moment it comes across that the 2D Stuart-Landau model and (predominantly linearly coupled) network versions dominate brain dynamic modelling, with a nod to the Wilson-Cowan model and a brief introduction to a next-generation mass model. Many other neural mass models are not mentioned, and there is a lack of discussion of popular phase reduction (and indeed phase-isostable) reductions that would give rise to phase oscillator networks of generalised Kuramoto type (the first type of model in the Table ``Summary of Neural Mass Model Equations''). Hopefully, the comments below will be useful in improving the exposition so that it can be a useful resource.}

\ourresponse{The structural changes prompted by this critique are: (i)~Section~2 has been retitled \emph{Linear Response and Phase-Only Limits} and given a rationale paragraph that justifies the linear-and-phase starting point on three concrete grounds (closed-form spectra, universal local description near a stable focus, graph-spectral / connectome-harmonic decomposition); (ii)~Section~3 has been compressed; (iii)~the roadmap table now lists the canonical whole-brain node models that are widely used in practice but not derived in detail in this review (Liley, Robinson, Wong--Wang/DMF, David--Friston, TVB); (iv)~explicit phase and phase--amplitude reductions of WILCO and NMM1 networks back to generalised Kuramoto-type descriptions are now cited at the points where the cross-walk requires them; (v)~Section~7 has been rebalanced to credit the Ott--Antonsen reductions of theta-neuron networks (Luke et~al.~2013; So et~al.~2014) that predate the QIF reduction of Montbri\'o, Paz\'o \& Roxin; and (vi)~the appendix structure has been pruned. The detailed responses are in the section-by-section notes below.
}

\section*{Sec 1 Introduction}

\reviewercomment{Given that this is a review, it would be useful to set the scene for why such a review is useful / necessary and how it moves on from existing or older reviews on similar topics. It might be useful to contrast with Ermentrout, G. B. Neural nets as spatio-temporal pattern forming systems. Rep. Prog. Phys. 61, 353-430 (1998). And P Ashwin et al. 2016 Mathematical frameworks for oscillatory network dynamics in neuroscience, Journal of Mathematical Neuroscience, 6:2 In the Introduction mass models are equated with rate models. The more common distinction is between point (mass) models and continuum (field) models. Both of these are often rate-based --- though not always. The statement `lacks a principled theoretical bridge' and surrounding text doesn't really make it clear whether this review will provide the bridge, whether it is an open challenge, or whether a new bridge is being proposed. The statement `Finally, we discuss the first neural mass model rigorously derived from first principles' is very misleading. The first do so this --- using a theta-neuron network --- were the following: So, P., Luke, T. B. \& Barreto, E. Networks of theta neurons with time-varying excitability: Macroscopic chaos, multistability, and final-state uncertainty. Phys D Nonlinear Phenom 267, 16-26 (2014). Luke, T. B., Barreto, E. \& So, P. Complete classification of the macroscopic behavior of a heterogeneous network of theta neurons. Neural Computation 25, 3207-3234 (2013). The link between this work and that on QIF has been reviewed in a number of places, e.g., [186].}

\ourresponse{The Introduction has been rewritten on five points. First, a short ``why now / what gap'' paragraph positions the review against Ermentrout~(1998) and Ashwin et~al.~(2016). Second, the mass--field distinction (lumped 0D vs.\ spatially continuous) replaces the mass--rate equation, and language equating ``neural mass'' with ``firing-rate'' has been removed; both formulations can be rate-, voltage-, or conductance-based. Third, the \emph{Rosetta Stone} metaphor is now framed as a practical cross-walk across modelling frameworks, not as a claim that no theoretical links exist. Fourth, the ``first rigorously derived'' claim is corrected: Luke et~al.~(2013) and So et~al.~(2014) are credited with the foundational Ott--Antonsen reductions of theta-neuron networks, and NMM2 is presented as a class of exact mean-field reductions exemplified by MPR rather than as a singular result. Fifth, a new paragraph within §1.1 \emph{Key Concepts} (inside the ``Neural mass vs.\ neural field'' item) makes the terminological scope explicit. Two usages of ``neural mass model'' coexist in the literature: a narrow one, in which the term names the Wilson--Cowan / Freeman / Lopes da Silva / Jansen--Rit / Wendling / LaNMM lineage and its exact next-generation extension by Montbri\'o, Paz\'o \& Roxin (current-based synapses, linear filters, and a transfer relation that may be a static sigmoid, a dynamically filtered sigmoid, or an exact dynamical $(r,v)$ relation derived from the spiking microscale), and a broader one in which ``neural mass'' simply means \emph{lumped} (as opposed to \emph{neural field}) and covers the conductance-based, propagation-speed, decision-circuit, and DCM families listed in block~(B) of the roadmap Table. We adopt the broader definition while stating clearly that the derivations concentrate on the rate/PSP lineage, where the chain of dynamical equivalences phase $\leftrightarrow$ damped $\leftrightarrow$ SL $\leftrightarrow$ WILCO $\leftrightarrow$ NMM1 $\leftrightarrow$ NMM2 is tightest; the broader families are credited and revisited in the synthesis but not re-derived in detail.}

\newpage
\section*{Section 1.1}

\reviewercomment{In Sec 1.1 (Key concepts) a case should be made as to why a neural mass model is represented by only a pair $(x,y)$, rather than a much higher dimensional system (point (a)). point (b) mixes periodic aperiodic, and chaotic behaviour without a clear distinction between them (as they are all described by `repetition'). point (d) needs the most unravelling. Chemical synaptic interactions are often `shunted', namely they depend on system state (not just input). Thus, a push could become a pull in models that respect conductance-based synapses. point (e) talks about weak forcing --- though one needs to also talk about strong attraction to cycle. In `Mathematical tools' the dimension of $x$ is not given, nor whether $f$ is smooth. `stability' should probably be `linear stability'.}

\ourresponse{Each item is addressed in the rewritten §1.1. We use the mass--field distinction; explicitly state that neural mass models in general have a state vector of model-dependent dimension, with the push--pull pair used only as a pedagogical core motif onto which the normal-form mappings hang; distinguish between self-sustained limit-cycle oscillations, damped oscillations, noise-driven quasi-cycles with spectral peaks, and chaotic dynamics with spectral structure; note that conductance-based synapses are state-dependent through the reversal potential, so the push--pull description is an effective one (typically a linearisation around an operating point); refine the forcing definition to require both weak forcing \emph{and} strong attraction to a stable limit cycle for the phase reduction to apply. The state-vector dimension and smoothness of $f$ are now stated where they are first used. Throughout, ``stability'' has been replaced by ``linear stability'' wherever the linearised criterion is the one being invoked.}

\section*{Table}

\reviewercomment{Looking at the Table ``Coupled Whole-Brain Models Summary of Neural Mass Model Equations'' one would expect to see here or elsewhere (e.g., in Sec 6 NMM1) models attributed to Liley, Robinson, Wong \& Wang, and others which all have distinct features not found in the models listed. e.g., that of Liley includes `shunts', the space-clamped version of Robinson emphasises axonal speeds, and Wong \& Wang is very popular. Dafilis, M. P., Frascoli, F., Cadusch, P. J. \& Liley, D. T. J. Four dimensional chaos and intermittency in a mesoscopic model of the electroencephalogram. Chaos 23, 023111 (2013). Breakspear, M. et al. A Unifying Explanation of Primary Generalized Seizures Through Nonlinear Brain Modeling and Bifurcation Analysis. Cereb. Cortex 16, 1296-1313 (2006). Wong, K.-F. \& Wang, X.-J. A Recurrent Network Mechanism of Time Integration in Perceptual Decisions. J Neurosci 26, 1314-1328 (2006). and used / reduced by Deco (Dynamic mean field (DMF) model) Deco G, Ponce-Alvarez A, Mantini D, Romani GL, Hagmann P, Corbetta M. Resting-state functional connectivity emerges from structurally and dynamically shaped slow linear fluctuations. J Neurosci. 2013;33(27):11239-11252 and see also Herzog, R. et al. Neural mass modeling for the masses: Democratizing access to whole-brain biophysical modeling with FastDMF. Netw. Neurosci. 8, 1590-1612 (2024). Other mass models are David, O. \& Friston, K. J. A neural mass model for MEG/EEG: coupling and neuronal dynamics. NeuroImage 20, 1743-1755 (2003). \dots\ or any implemented in The Virtual Brain (thevirtualbrain.org). \dots\ and for applications see e.g., Bastiaens, S. P., Momi, D. \& Griffiths, J. D. A comprehensive investigation of intracortical and corticothalamic models of the alpha rhythm. PLOS Comput. Biol. 21, e1012926 (2025). Theoretical Neuroscience: Understanding Cognition, X.-J. Wang, CRC Press, 2025.}

\ourresponse{The roadmap table has been revised in two ways. First, the Kuramoto / damped-oscillator / Stuart--Landau entries are now explicitly labelled as a normal-form ladder used as a pedagogical entry point, not as a claim of dominance. Second, a new block lists canonical whole-brain node models that are widely used but are not derived in detail in this review: conductance-based / shunting Liley-type masses (Dafilis et~al.~2013), Robinson-style corticothalamic / propagation-speed models (Robinson 2001; Rennie et~al.~2002), Wong--Wang and dynamic mean-field (DMF) nodes (Wong \& Wang 2006; Deco et~al.~2013), and David--Friston / DCM neural masses (David \& Friston 2003), with a one-sentence pointer to The Virtual Brain (Sanz Le\'on et~al.~2013, 2015). This places the cross-walk inside the broader landscape rather than letting the oscillator normal forms appear to occupy it.}

\section*{Sec 2 Harmonic Oscillator}

\reviewercomment{It is not really clear that a physics audience needs a section with this title. Certainly, the rationale for treating linearly coupled harmonic oscillators needs to be given. Moreover, in a physical or neurobiological sense what does it mean to linearly couple in the complex state variable? Eqn (2.17) should be complemented with the dynamics for $r_i$. Provide a reference for the Kuramoto-Sakaguchi structure after eqn (2.17). For eqn (2.20) there should be some neurobiological discussion of the meaning of $C_{ij}$ and $\tau_{ij}$ --- structural connectome and delays from white matter tracts, all potentially available from human connectome repositories. In Sec 2.1.5 Applications, alpha frequency is mentioned, though never defined or discussed in a neuroscience context. Indeed, what about the common frequency bands used in EEG --- never really discussed anywhere. Building up from SL to linear coupled SL networks is fine --- but why? What is the link to either neuroscience or the key mathematical modelling idea? How does one numerically evolve (2.29) --- it is both stochastic and delayed. Sec 2.1.5 Applications mentions several technical terms, e.g., `graph-spectral neural field' and `Laplacian', without technical qualification (at odds with the slow build-up of previous sections) and then makes a jump to `pharmacology' that is very hard to fathom or understand without context, qualification, or detail.}

\ourresponse{The section has been retitled \emph{Linear Response and Phase-Only Limits} and given a methodological rationale up front: linearised damped oscillator networks admit closed-form covariances, lagged covariances, and power spectra (and so analytic predictions for FC and MEG/EEG PSDs); every nonlinear neural mass model introduced later (SL, WILCO, NMM1, NMM2) reduces near a stable focus to a linear damped oscillator, so the section supplies the universal local description; and the linear network on a connectome is the natural setting for graph-spectral and connectome-harmonic decompositions. The phase reduction is now stated explicitly as a modelling choice valid under weak forcing and strong attraction to a limit cycle. The amplitude equation is written alongside the phase equation, and the phase-only model is derived as the slow-manifold reduction obtained when amplitudes relax to a fixed spatial profile.

The Sakaguchi--Kuramoto form is named at first appearance, with the offset $\iota_{ij}$ identified as the small-delay reduction $\iota_{ij}\approx\bar\omega\,\tau_{ij}$ at a carrier frequency. Connectivity $C_{ij}$ and delay $\tau_{ij}=d_{ij}/v$ are now identified with the structural connectome and white-matter conduction delays ($v\in[5,10]$~m/s). EEG/MEG bands are defined in a short box (Box~2.1) with their conventional ranges and dominant generators; subsequent uses of ``alpha'' refer back to that box. Graph-spectral terminology is introduced in the rationale paragraph.

The numerics correction is the most substantive change. We no longer claim that Euler--Maruyama handles delays. The dynamics with non-zero delays and stochastic forcing are stochastic delay differential equations; we recommend method-of-steps DDE solvers in the deterministic limit, frame fixed-step Euler--Maruyama on the delayed system as a controlled approximation, and specify the noise discretisation, the interpolation scheme on the delay buffer, the dual step-size bound $\Delta t\le\min\{1/(50\,f_{\max}),\,\min_{ij}\tau_{ij}/10\}$, and a convergence check by step halving. Reproducibility reporting (solver, version, scheme, $\Delta t$, simulated time, transient length, RNG seed) is now stated explicitly. The same correction has been applied to the analogous boxes in §§3, 6, 7. The pharmacology paragraph has been rewritten so that stimulation and pharmacological modulation appear as two examples of the same mechanism (designed perturbations of effective gain and damping near a stable focus) rather than as an abrupt jump.}

\section*{Sec 3 Stuart Landau Oscillator}

\reviewercomment{Again, it is unclear why a physics audience needs such a review (or historical perspective). Much of this section is well covered by [51,71] --- what is the need to repeat it? This Sec could be sacrificed / shortened at the expense of the topics the audience is likely to be less familiar with, namely large-scale brain modelling and the plethora of neural mass models in this space. Homeostatic mechanisms are mentioned, and relevant papers are Vogels T.P., Sprekeler H., Zenke F., Clopath C., Gerstner W. Inhibitory plasticity balances excitation and inhibition in sensory pathways and memory networks. Science. 2011;334(6062):1569-1573 Nicola, W., Hellyer, P. J., Campbell, S. A. \& Clopath, C. Chaos in homeostatically regulated neural systems. Chaos Interdiscip J Nonlinear Sci 28, 083104 (2018). Al-Darabsah, I., Campbell, S. A. \& Rahman, B. Distributed Delay and Desynchronization in a Neural Mass Model. SIAM J. Appl. Dyn. Syst. 23, 3013-3051 (2024). A nice use of noisy SL in neuroscience is the following and should be mentioned: Freyer, F., Roberts, J. A., Ritter, P. \& Breakspear, M. A Canonical Model of Multistability and Scale-Invariance in Biological Systems. PLoS Comput. Biol. 8, e1002634 (2012). The complexified Kuramoto model might also be worth mentioning in this section: Budzinski, R. C. et al. Analytical prediction of specific spatiotemporal patterns in nonlinear oscillator networks with distance-dependent time delays. Phys. Rev. Res. 5, 013159 (2023).}

\ourresponse{Section~3 has been compressed. The Landau / Stuart historical material is now a brief acknowledgement plus a pointer to the standard references; the polar / Cartesian / complex equivalence is retained briefly in the main text for the cross-walk to Wilson--Cowan, NMM1, and NMM2 — where each form recurs naturally — and the longer textbook-style derivation is given in Appendix~C.

We retain only what is needed for the cross-walk: SL as the Hopf normal form (with the bifurcation parameter as the local working point biophysical neural mass models inherit near a Hopf), the shear $\beta$ as the parameter that distinguishes SL from the linear damped oscillator, and the noise-near-criticality regime in which SL networks are typically operated.

The reviewer's specific suggestions are taken on board. Freyer et~al.~(2012) is now the empirical anchor for the noise-driven multistability discussion: scale-invariant alpha-rhythm bistability under noise near a subcritical Hopf in an SL-type model. Budzinski et~al.~(2023) is cited at the SL-network polar reduction, where their analytical predictions for spatiotemporal patterns under distance-dependent delays are directly relevant to whole-brain SL models on a connectome. The previous mention of ``homeostatic mechanisms'' as a hand-wave for nonlinear saturation is replaced by explicit citations to Vogels et~al.~(2011) on inhibitory plasticity for E/I balance, Nicola et~al.~(2018) on chaos in homeostatically regulated systems, and Al-Darabsah, Campbell \& Rahman (2024) on distributed-delay neural mass models.}

\section*{Sec 4 Interlude}

\reviewercomment{Although conductance changes can nicely be modelled with temporal filters, synaptic currents are shunted. Namely, synapses have reversal potentials and interactions are state dependent. These are not mentioned or covered. In Sec 4.1 citations to Jansen-Rit, Wendling, or LaNMM are missing.}

\ourresponse{A new subsection §4.2, \emph{Current-based vs.\ conductance-based synapses}, has been added immediately after the operator formalism (the previous §4.2 \emph{Transfer functional} is renumbered §4.3). The new subsection states that the linear-filter description retained throughout this review is a current-based idealisation and gives the canonical conductance-based form $I_{\rm syn}(t)=g(t)\bigl[E_{\rm rev}-V(t)\bigr]$, making the shunting effect concrete: as $V(t)\to E_{\rm rev}$ the driving force vanishes and the synapse delivers no current despite an elevated conductance, and if $V$ overshoots $E_{\rm rev}$ the effective sign of the interaction reverses. We retain the current-based form for the cross-walk because the rate-based and exact-mean-field models in §§5--7 are themselves current-based, and because the conductance-based interaction linearises to a current-based effective coupling near any stable operating point. The conductance-based / shunting Liley-type mass listed in block~(B) of the roadmap table is the natural starting point for readers who require the full $V$-dependent treatment, including the NMDA voltage dependence already noted at the end of §4.1. The previous unsupported reference in §4.1 now carries explicit citations to Jansen \& Rit, the Wendling extensions, and the laminar / LaNMM work.}

\section*{Sec 5 Wilson-Cowan}

\reviewercomment{The statement that WILCO is SL is too bold for me. I would only expect this to be true in the nbhd of a Hopf bifurcation. Sec 5.3 should emphasise that this model (eqn 5.5) has been studied by many others and give some citations and summaries. In Sec 5.4 I would argue that the WILCO model does not `trade spiking detail' as it is not derived from a spiking model. For reductions of WILCO networks to phase oscillator networks see e.g., Hlinka J, Coombes S. Using computational models to relate structural and functional brain connectivity. Eur J Neurosci. 2012;36:2137-45. Daffertshofer A, van Wijk BCM. On the influence of amplitude on the connectivity between phases. Front Neuroinform. 2011;5:6.}

\ourresponse{The reviewer is right on both substantive points. The previous claim that ``centre-manifold theory confirms that WILCO is SL in disguise'' is replaced by ``near a Hopf bifurcation, WILCO reduces on its centre manifold to the SL normal form''; a sentence has been added making explicit that this is a strictly local equivalence and that away from the bifurcation WILCO supports saturation regimes, multistability, and pattern-formation regimes that are absent from the SL normal form. The full local reduction remains in Appendix~G (relabelled \Verb|app:WILCO2SL|).

The lineage statement is corrected. WILCO is now described as a phenomenological rate model formulated at the population level, not as a reduction of spiking dynamics. §5.3 \emph{Coupling} now opens with a paragraph placing the network form within the established literature, citing Borisyuk \& Kirillov, Hoppensteadt \& Izhikevich, and the more recent multistability / metastability and bifurcation-theoretic studies. A new paragraph at the end of §5.3 introduces phase and phase--amplitude reductions of WILCO networks as the explicit bridge back to generalised Kuramoto-type descriptions, citing Hlinka \& Coombes (2012) and Daffertshofer \& van Wijk (2011); this closes the WILCO\,$\to$\,Kuramoto leg of the cross-walk that the reviewer rightly noted was missing.}

\section*{Sec 6 NMM1}

\reviewercomment{The Barkhausen stability criterion is deferred to an appendix (a pointer to J3 should be given). It would have been more useful to see this discussed in context and in appropriate detail --- especially as it is important to the push/pull narrative. In eqn (6.5) why are there only $\sim\!N$ delays for $\sim\!N^{2}$ connections? The laminar model is not well described and essentially devolved to Fig 6.1. For the reduction of Jansen-Rit, see e.g., M Forrester et al. 2020 The role of node dynamics in shaping emergent functional connectivity patterns in the brain, Network Neuroscience, 4:2, 467-483.}

\ourresponse{A condensed half-page Barkhausen exposition now sits in the main text immediately after the second-order operator: the loop-gain and loop-phase conditions in their classical form, an explanation of why the second-order synaptic filter (with distinct rise and decay) supplies the frequency-dependent phase shift required to close the loop, and a statement that the push--pull E$\to$I$\to$E motif is the minimal architecture that satisfies the criterion in NMM1. The longer pedagogical exposition with the algebra remains in Appendix~J.3 (now labelled \Verb|app:barkhausen|).

The delay indexing in Eqs.~6.5--6.6 was an error; both equations have been corrected to the edge-indexed form $\tau_{ij}$ ($\sim\!N^{2}$ delays), matching the convention used throughout the rest of the manuscript and in the roadmap table.

A sentence at the end of §6.1 now points to Forrester et~al.~(2020) for the explicit phase reduction of Jansen--Rit-type NMM1 networks and its consequences for the emergent functional connectivity patterns observed in whole-brain simulations, completing the NMM1\,$\to$\,phase-oscillator leg of the cross-walk.

The LaNMM paragraph in §6.2 has been expanded. It now describes the two coupled generators (a Jansen--Rit-like deep slow generator producing alpha/theta, and a PING-like superficial fast generator producing gamma) and their cross-frequency coupling; the laminar projection structure and the volume-conduction physics that link the model to depth-LFP and CSD observables; the resulting depth-dependent polarity and phase relations; and the empirical anchors --- laminar spectral peaks and CSD sinks/sources in primate and rodent cortex (Mendoza-Halliday et~al.~2024; Buffalo et~al.~2011; Spaak et~al.~2012; Sotero et~al.~2015), Alzheimer's-related alpha slowing and gamma loss, and tES/tACS coupling.}

\section*{Sec 7 NMM2}

\reviewercomment{It is stated that the earlier models in NMM1 are partly derived from 1st principles and later they are referred to as phenomenological --- so which is it? This section should be re-balanced to cover the earlier work of So, Luke, \& Barreto that predates MPR. After eqn (7.6) --- note that a Lorentzian distribution does not have a mean (so change mean to center). All of the eqns should be given in (7.8) (namely those missing in \dots). Moreover, others have considered such networks and also with delays, e.g., [200] with second order synapses and first order synapses in: Rabuffo, G., Fousek, J., Bernard, C. \& Jirsa, V. Neuronal Cascades Shape Whole-Brain Functional Dynamics at Rest. Eneuro 8, ENEURO.0283-21.2021 (2021). Other relevant work is Mayora-Cebollero, A., Barrio, R., Li, L., Mayora-Cebollero, C. \& P\'{e}rez, L. Dynamics of coupled neural populations: The role of synaptic dynamics. Chaos: Interdiscip. J. Nonlinear Sci. 35, 063140 (2025). Another type of NMM2 model is Nicola, W. \& Campbell, S. A. Mean-field models for heterogeneous networks of two-dimensional integrate and fire neurons. Front Comput Neurosc 7, 184 (2013).}

\ourresponse{The lineage is rebalanced. The §7 introduction now credits Luke, Barreto \& So (2013) and So, Luke \& Barreto (2014) for the exact macroscopic descriptions of theta-neuron networks via the Ott--Antonsen ansatz before introducing MPR (Montbri\'o, Paz\'o \& Roxin~2015) as the parallel reduction for QIF networks. Because QIF and theta-neuron networks are related by the standard quadratic-to-theta change of variables $V=\tan(\theta/2)$, the two reductions describe the same macroscopic dynamics in different coordinates; NMM2 inherits from both lines.

The ``phenomenological vs.\ first-principles'' inconsistency is removed. The text now states cleanly that in NMM1 the synaptic filter $\hat L$ is biophysically grounded (mass-action kinetics, PSP impulse response) while the static wave-to-pulse sigmoid is phenomenological (Freeman); NMM2 retains the biophysical synaptic filter and replaces the phenomenological sigmoid with a dynamical $(r,v)$ relation derived exactly from the QIF microscale via the Ott--Antonsen manifold.

The Lorentzian terminology is corrected. After Eq.~7.6 the parameters $\bar\eta$ and $\Delta$ are described as ``the centre and half-width-at-half-maximum (HWHM) of the Lorentzian (Cauchy) distribution.'' A parenthetical notes that the Lorentzian admits no moments of order $\ge 1$ (the integral defining the mean is divergent), so $\bar\eta$ is a location parameter rather than an expectation; under the Cauchy principal value, $\bar\eta$ formally equals the median. The corresponding NMM2 box description has been synchronised (``Centre (location parameter) and HWHM'').

The boxed E--I network (\Verb|eq:NMM2-EI|) now displays the explicit six-equation system $(\dot r_x,\dot v_x,s_x;\,\dot r_y,\dot v_y,s_y)$ in place of the previous static-rate form that used $\hat\Phi$ and hid the voltage dynamics. The accompanying text states the sign convention explicitly --- all coupling magnitudes $C_{\alpha\beta}\ge 0$, with $+$ for excitatory drive ($C_{xx}$, $C_{yx}$, long-range $C_{ij}$) and $-$ for inhibitory drive ($C_{yy}$, $C_{xy}$) in the voltage equations --- so that the swap-and-flip rule from the excitatory ($x$) to the inhibitory ($y$) population is read directly off the equations.

A short paragraph at the end of the §7 introductory material situates NMM2 as one instance within a broader family of exact mean-field reductions: Rabuffo et~al.~(2021) on whole-brain functional dynamics from NMM2-class networks with delays; Mayora-Cebollero et~al.~(2025) on the role of synaptic dynamics in coupled MPR populations; and Nicola \& Campbell (2013) on exact mean-field reductions for heterogeneous networks of two-dimensional integrate-and-fire neurons (an alternative beyond the QIF/theta line).

The numerics correction extends to NMM2. The ``How to Use It'' box for the NMM2 multi-population network previously read ``integrate with RK / Euler--Maruyama, $\Delta t\le 1/(200\,f_{\max})$'' without acknowledging that the equations with delayed coupling and stochastic forcing are SDDEs. The box has been rewritten in parallel with the §2 correction.}

\section*{Sec 8 Summary and Outlook}

\reviewercomment{How are we to see the synthesis across the models discussed --- this Section should have a clearer message in this regard. Others have written good perspectives on neural masses that may be of inspiration: Griffiths, J. D., Bastiaens, S. P. \& Kaboodvand, N. Computational Modelling of the Brain, Modelling Approaches to Cells, Circuits and Networks. Adv. Exp. Med. Biol. 1359, 313-355 (2021).}

\ourresponse{The closing has been rewritten so that the synthesis is explicit. A new opening paragraph states the unifying message: across all formalisms (linear damped, SL, WILCO, NMM1, NMM2), a single core motif recurs --- a push--pull pair of effective degrees of freedom, coupled through linear filters and a nonlinear transfer function, embedded in a network through structured coupling and delays. In the linear limit this core reduces to damped resonators or complex Ornstein--Uhlenbeck processes with closed-form covariances and spectra; near a Hopf bifurcation it reduces to the SL normal form; WILCO and NMM1 add biophysical levers without losing this dynamical structure; and NMM2 shows that, for QIF networks at least, the macroscopic descriptions can be derived exactly rather than postulated.

A paragraph now positions the present derivation--representation cross-walk relative to the spatial-scale organisation of Griffiths, Bastiaens \& Kaboodvand~(2021), recommending the two perspectives as complementary entry points to the field. The outlook is structured along five axes: extending exact mean-field reductions beyond QIF; more realistic local circuits (multiple inhibitory subtypes, chemical and electrical synapses, short-term plasticity, intrinsic adaptation); whole-brain heterogeneity and laminar-specific directed coupling, motivated by predictive-processing accounts and cytoarchitectonic / receptor / gene-expression gradients; neuromodulatory systems as a low-dimensional control axis; and connections to information-theoretic and Bayesian inference frameworks.}

\section*{Appendices}

\reviewercomment{I find many of these to be unnecessary for the main body. A: Going over names/terminology doesn't add anything. If it is not spelled out in the main body then it should be. B: Plane to polar coords should be within the grasp of the intended audience. C: Standard textbook material --- pointers to the literature can be given instead. The bifurcation diagrams for the various mass models are useful and could be moved to the main body (and possibly some two parameter diagrams prepared too). D: Standard textbook material. E: OK --- but is it necessary? It does not appear to be used in the main body. F: Query: why is it OK to drop some of the terms in the expansion of $(\Delta h_x)^2$? e.g., aren't $(x-x_0)^2$ and $(x-x_0)\,\delta I_x$ of the same order? If not why treat $(u_x, u_y)$ and $(\delta I_x, \delta I_y)$ to be of different orders? G: OK --- but expected as normal form for Hopf. H: Standard textbook material. I: Koopman operators not even mentioned in the main body. What is the need for this appendix? The boxes throughout the review (including when / how to use) are very repetitive and not always informative. e.g. why only track Kuramoto and not Kuramoto-Daido parameters in `Coupled Undamped Harmonic Oscillators', and numerical schemes are often not appropriate (e.g., Euler-Maruyama does not handle delays).}

\ourresponse{We accept the underlying principle that appendices should support the main narrative or be removed. Appendices~A, B, D, and~H have been compressed: only material actively used in the main text is retained, and where the content is genuinely standard we now point to standard references (Strogatz; Kuramoto; Coombes \& Bressloff; Ermentrout \& Terman) rather than reproduce the derivations.

The most informative one-parameter bifurcation diagrams of the WILCO and NMM1 single-population models (WC, Jansen--Rit, Wendling, LaNMM) have been moved into the main body --- specifically into Sections~5 and~6, where they now anchor the discussion of where each model lives in parameter space. The remaining bifurcation diagrams (Stuart--Landau, ING, PING) are kept in the appendix as a complementary gallery, with a breadcrumb pointing to the new main-body locations. We have not added two-parameter diagrams in this revision; the existing one-parameter set covers the qualitative parameter-space organisation needed for the cross-walk.

Appendix~E (\emph{Linear coupled complex oscillators}) is retained and is now operationally cited from §2.1.2, where the polar reduction in the main text needs the full eigenmode decomposition and the U(1)-symmetry argument that Appendix~E supplies. Appendix~G is retained as the canonical centre-manifold reduction underpinning the local equivalence statement in §5; the cross-reference from §5 has been corrected.

Appendix~F is the more delicate point. The reviewer is right that retaining the cross-terms $S''(h_{x,0})\,w_{xx}\,\delta I_x\,u_x$ and $S''(h_{x,0})\,\delta I_x\,w_{xy}\,u_y$ alongside the dropped $u_x^2$ and $u_x u_y$ terms is inconsistent without an explicit scaling argument. We have chosen Option~F3: state deviations $u_x,u_y=\mathcal{O}(\varepsilon)$ and externally controlled input modulation $\delta I=\mathcal{O}(\varepsilon^{1/2})$, so that mixed terms $u\cdot\delta I=\mathcal{O}(\varepsilon^{3/2})$ are retained while pure-state quadratic terms $u^2=\mathcal{O}(\varepsilon^2)$ are dropped. This is the natural ordering when $\delta I$ is an experimentally controlled drive and $u$ is the system's response, and it grounds the FM-modulation interpretation in the latter half of the appendix.

Appendix~I (Koopman) was thinly motivated in the original. %We have not added a phase / isostable bridge in the main body, since that would inflate the review beyond its scope; instead the appendix has been pruned to a half-page summary that grounds the algorithmic-information / compression framing of ``what counts as an oscillation'' already present in §1.1, with two sentences linking Koopman eigenfunctions to the phase--amplitude descriptions used in §§2--3.
Koopman material has been folded into Appendix~H (Oscillations, Topology and Simplicity), where it now serves as the operator-theoretic component of the algorithmic-information framing rather than a standalone appendix.

The ``How to use it'' boxes have been corrected in two ways. First, the Euler--Maruyama-with-delays language is replaced by the SDDE / method-of-steps formulation across the relevant boxes (§§2.1, 3, 6, 7), with the full reproducibility-reporting list (solver, version, discretisation scheme, $\Delta t$, simulated time, transient length, RNG seed) consolidated in the §2.1 box and cross-referenced from the others. Second, the repeated boilerplate has been trimmed while keeping each box self-contained for the model it documents.

On the reviewer's specific note that the Coupled Undamped Harmonic Oscillator box tracks only the standard Kuramoto order parameter $Z(t)=\tfrac{1}{N}\sum_j e^{i\theta_j(t)}$ rather than the full Kuramoto--Daido hierarchy: we retain $Z(t)$ as the box readout because it is the order parameter for which the Sakaguchi--Kuramoto and graph-spectral results derived in §2 are expressed. Higher-order Kuramoto--Daido coefficients $Z_n=\langle e^{i n\theta}\rangle$ can be added as box readouts if the reviewer regards this as essential.}

\section*{Typos}

\reviewercomment{Before eqn (4.20): $<\!<$ should be $\ll$ The References could be much better polished. e.g., names in paper titles need capitalisation (kuramoto $\to$ Kuramoto), and Ref [1] and [249] are the same.}

\ourresponse{The $<\!<$ before Eq.~(4.20) is now $\ll$ and the rest of the manuscript has been checked for the same artefact. The bibliography has been audited and proper nouns are now braced (Kuramoto, Hopf, Stuart, Landau, Wilson--Cowan, Jansen--Rit, Wendling, Ott--Antonsen, Hodgkin--Huxley, Hilbert, Lorentzian, Cauchy, Bayesian, etc.) so that they retain their capitalisation under the journal's bibliography style. Ref~[1] (Ashwin, Coombes \& Nicks 2016) and Ref~[249] (Kringelbach et~al.~2020) are distinct entries in the current version; if the reviewer was working from an earlier file in which they were duplicated, the de-duplication is now done as part of the bibliography polish.}

\end{document}
